\NeedsTeXFormat{LaTeX2e}
\documentclass[a4paper,12pt,
headsepline,           % Linie zw. Kopfzeile und Text
oneside,               % einseitig
pointlessnumbers,      % keine Punkte nach den letzten Ziffern in Überschriften
bibtotoc,              % LV im IV
%DIV=15,               % Satzspiegel auf 15er Raster, schmalere Ränder   
%BCOR15mm               % Bindekorrektur
%,draft
]{scrartcl}

\usepackage{amsmath}
\usepackage{amsfonts}
\usepackage{amssymb}
\usepackage{enumitem}
\usepackage[utf8]{inputenc} % this is needed for umlauts
\usepackage[ngerman]{babel} % this is needed for umlauts
\usepackage[T1]{fontenc} 
\usepackage{commath}
\usepackage{xcolor}
\usepackage{booktabs}
\usepackage{float}
\usepackage{tikz-timing}
\usepackage{tikz}
\usepackage{multirow}
\usepackage[final]{pdfpages}
\usepackage{blindtext}
\usepackage[scaled]{helvet}
\usepackage{hyperref}

\usetikzlibrary{calc,shapes.multipart,chains,arrows}

\KOMAoptions{DIV=last} % Neuberechnung Satzspiegel nach Laden von Paket helvet

\usepackage{scrpage2}
\pagestyle{useheadings}

\renewcommand{\familydefault}{\sfdefault} 

\setlength{\parindent}{0pt}   % kein linker Einzug der ersten Absatzzeile
\setlength{\parskip}{1.4ex plus 0.35ex minus 0.3ex} % Absatzabstand, leicht variabel

\newcommand{\fullname}{Gruppe 10}
\newcommand{\titel}{Softwaregrundprojekt Meilenstein 1}
\newcommand{\jahr}{2018}
\newcommand{\dozent}{Florian Ege}
\newcommand{\betreuer}{Stefanos Mytilineos}
\newcommand{\fakultaet}{Ingenieurwissenschaften, Informatik und\\Psychologie}
\newcommand{\institut}{Institut für Softwaretechnik und Programmiersprachen}

\pdfinfo{
    /Author (\fullname)
    /Title (\titel)
    /Producer     (pdfeTex 3.14159-1.30.6-2.2)
    /Keywords ()
}

\hypersetup{
    pdftitle=\titel,
    pdfauthor=\fullname,
    pdfsubject={Softwaregrundprojekt-Abgabe},
    pdfproducer={pdfeTex 3.14159-1.30.6-2.2},
    colorlinks=false,
    pdfborder=0 0 0	% keine Box um die Links!
}

% Trennungsregeln
\hyphenation{Sil-ben-trenn-ung}


\newcommand{\begriff}[7] {
    \begin{table}[H]
        \centering
        \begin{tabular}{|p{4.5cm}|p{10cm}|}
            %\hline
            %\toprule \\
            \hline
            \textbf{Begriff} & \textbf{#1} \\ \hline
            %\midrule \\
            \textbf{Beschreibung} & #2 \\ \hline
            %\midrule
            \textbf{Ist-ein} & #3 \\ \hline
            %\midrule
            \textbf{Kann-sein} & #4 \\ \hline
            %\midrule
            \textbf{Aspekt} & #5 \\ \hline
            %\midrule
            \textbf{Bemerkung} & #6 \\ \hline
            %\midrule
            \textbf{Beispiel} & #7 \\ %\hline
            %\bottomrule
            \hline
        \end{tabular}
    \end{table}
}

%\title{Softwaregrundprojekt}
%\author{Gruppe 10}

\newcommand{\anf}[7] {
    \begin{table}[H]
        \centering
        \begin{tabular}{|p{3.2cm}|p{9.8cm}|}
            \hline 
            \textbf{ID:} & \textbf{#1} \\ \hline
            \textbf{Titel:} & #2 \\ \hline
            \textbf{Beschreibung:} & #3 \\ \hline
            \textbf{Begründung:} & #4 \\ \hline
            \textbf{Abhängigkeiten:} & #5 \\ \hline
            \textbf{Priorität:} & #6 \\ \hline
            \textbf{Akteure:} & #7 \\ \hline
        \end{tabular}
    \end{table}
}

\newcounter{fanfCount}
\newcommand{\fanf}[6] {
    \stepcounter{fanfCount}
    \anf{FA\thefanfCount}{#1}{#2}{#3}{#4}{#5}{#6}
}
\newcounter{qanfCount}
\newcommand{\qanf}[6] {
    \stepcounter{qanfCount}
    \anf{QA\theqanfCount}{#1}{#2}{#3}{#4}{#5}{#6}
}

\newcommand{\akt}[4] {
    \begin{table}[H]
        \centering
        \begin{tabular}{|p{3cm}|p{10cm}|}
            \hline
            \textbf{ID:} & \textbf{#1} \\ \hline
            \textbf{Titel:} & #2 \\ \hline
            \textbf{Beschreibung:} & #3 \\ \hline
            \textbf{Rolle:} & #4 \\ \hline
        \end{tabular}
    \end{table}
}

\newcounter{faktCount}
\newcommand{\fakt}[3] {
    \stepcounter{faktCount}
    \akt{AKT\thefaktCount}{#1}{#2}{#3}
}

\begin{document}
    \thispagestyle{empty}
    \begin{addmargin*}[4mm]{-10mm}

        \includegraphics[height=1.8cm]{images/unilogo_bild}
        \hfill
        \includegraphics[height=1.8cm]{images/unilogo_wort}\\[1em]

        {\footnotesize
        %{\bfseries Universität Ulm} \textbar ~89069 Ulm \textbar ~Germany
        \hspace*{115mm}\parbox[t]{35mm}{\bfseries Fakultät für\\
        \fakultaet\\
        \mdseries \institut}\\[2cm]

        \parbox{140mm}{\bfseries \LARGE \titel}\\[2.5em]
        {\footnotesize Softwaregrundprojekt an der Universität Ulm}\\[3em]

        {\footnotesize \bfseries Vorgelegt von:}\\
        {\footnotesize \fullname\\}\\ [1em]
        {\footnotesize \bfseries Dozent:}\\
        {\footnotesize \dozent\\}\\[1em]
        {\footnotesize \bfseries Betreuer:}\\
        {\footnotesize \betreuer}\\ [1em]
        {\footnotesize \jahr}
        }
    \end{addmargin*}
    \pagebreak
    \tableofcontents
    \pagebreak

    \section{Kontextanalyse}
    \subsection{Einleitung}
Bei dem Projekt handelt es sich um die Konzeption und Implementierung eines online Multiplayer-Spiels aus der Welt von Harry Potter – genauer: \textit{Fantastic Feasts}. Es ist eine rundenbasierte Form des bekannten Spiels Quidditch.


Im Mittelpunkt des Projekts steht das Erlernen von Fähigkeiten im Umgang mit einem größeren Softwareprojekt. Es werden keine kommerziellen Ziele verfolgt.

Der Auftraggeber – im weitesten Sinne die Universität Ulm – verfolgt das Ziel, den Studenten Fähigkeiten zu vermitteln und sie anschließend nach genau definierten Maßstäben zu bewerten. Zu diesen Fähigkeiten gehört folgendes:

Zunächst einmal stehen Planen, Formulieren von Anforderungen und Modellierung von Software an. Es folgt die Auseinandersetzung mit verschiedenen Plattformen und Technologien auf die für die Implementierung zurückgegriffen werden soll. Gleichzeitig wird das Ziel verfolgt, übergeordnete Fähigkeiten zur Qualitätssicherung, zur Versionenverwaltung oder zu agilen Entwicklungsprozessen im Team zu erwerben. 
Erst dann kommen praktische Programmierfähigkeiten zum tragen. Auch hier ist es das Ziel, diese auszubauen.

Die Studierenden – in diesem Fall 6 Studenten der Informationssystemtechnik – verfolgen das Ziel, das Projekt nach den Anforderungen im Lastenheft erfolgreich umzusetzen und die Abnahmeprüfung zu bestehen.





\subsection{Motivation}
Die Motivation für das Projekt lässt sich – wie bei den Zielen – in die der Universität und die der Studenten aufteilen.

Die Universität gibt die Inhalte vor. Diese sollen von den Studenten bestmöglich erlernt werden, da es Teil ihrer Ausbildung darstellt. Die Universität will somit ihrem Auftrag der Lehre gerecht werden.
 

Für die Studenten ist die Motivation der Erwerb und Ausbau der oben genannten Fähigkeiten und im weitesten Sinne eine erfolgreiche Ausbildung in ihrem Fach. Zusätzlich soll ein Spiel entwickelt werden, das funktioniert und Spaß macht.
 

\subsection{Vision}
Das fertige System soll folgendermaßen aufgebaut sein: Einer Client-Server-Architektur folgend kommunizieren ein oder mehrere Clients mit dem Server, auf dem die Spiellogik läuft. Die Spieler haben client-seitig eine ansprechende und lebendige GUI, über die sie \textit{Fantastic Feasts} spielen, eine Partie als Zuschauer verfolgen, Charaktere und Ausrüstung zu Teams mitsamt Farben und Logo zusammen stellen können und die Möglichkeit haben, Partien zu konfigurieren. Begleitet wird die visuelle Darstellung von Soundeffekten und einer thematisch Ansprechenden Spielmusik.

Im durch und durch taktischen Spiel mit zwei sich genüberstehenden Quidditch-Teams können die Spieler Runde für Runde Spielfiguren auf Besen über das Spielfeld jagen lassen, Punkte erzielen, den Gegner sabotieren und Publikumseffekte zu ihrem Vorteil einsetzen. Doch selbst bei noch so guter Taktik kann ihnen der Zufall einen Strich durch die Rechnung machen, da nicht immer alles so eintritt, wie es sich die einzelnen Spieler vielleicht erhofft haben.

Was die Spielmodi betrifft, ist das Kern-Szenario das Multiplayer-Spiel. Hier entsteht durch den Wettstreit zweier Spieler die größte Spannung. Doch um das Spiel auch alleine spielbar zu machen, existiert ein Singleplayer-Modus. Eine ausgefeilte KI mit voraussichtlich mehreren Schwierigkeitsstufen stellt für Einzelspieler eine spannende Herausforderung dar. Wer sich noch auf den großen Wettkampf vorbereitet, die Taktik andere Spieler erlernen will oder einfach Spaß am Zuschauen hat, kann sich im Zuschauer-Modus in andere Multiplayer-Partien einklinken. Damit bleibt einer breiten Zielgruppe an Spielern kaum etwas zu wünschen übrig.

\subsection{Projektkontext}
Auftraggeber des Projektes ist die Servicegruppe Informatik der Universität Ulm, die das Modul Softwaregrundprojekt veranstaltet. Der Tutor Stefanos Mytilineos vertritt den Auftraggeber während der Projektlaufzeit, in höherer Instanz ist Florian Ege verantwortlich. Als weiterer Stakeholder tritt das Team auf, das das Projekt letztendlich entwickelt. Es besteht aus sechs Studenten der Informationssystemtechnik: Tarik Enderes, Tim Luchterhand, Jonas Merkle, Paul Nykiel, Björn Petersen und Michael von Hohnhorst. Diese bearbeiten das Projekt nicht in Vollzeit, da sie parallel den weiteren Verlauf ihres Studiums verfolgen. 

Indirekt am Entwicklungsprozess beteiligt ist das Standardisierungskomitee, in das auch aus diesem Team ein Vertreter geschickt wird. Dort werden alle nötigen Protokolle und Schnittstellen definiert, die bei der Entwicklung von \textit{Fantastic Feasts} benötigt werden. Weitere Stakeholder, die jedoch erst später in Erscheinung treten, sind andere Teams, die gegebenenfalls auf einen Komponenten dieses Projektes angewiesen sind. Sie werden auf einer bevorstehenden Messe zu potenziellen Kunden. In ihrem Interesse liegt eine saubere Implementierung bei gleichzeitig guter Dokumentation des Komponenten. Zu guter letzt muss dieses Projekt sowie die Einzelleistung eines jeden Team-Mitgliedes die Prüfer in der Abnahmeprüfung überzeugen – sie stellen so gesehen die wichtigsten Kunden dar.

Da bei der Implementierung auf den agilen Entwicklungsprozess Scrum zurückgegriffen wird, sollen auch hier die Rollen kurz benannt werden. In dieser Phase übernimmt der oben genannte Tutor die Rolle des Product Owners. Der ScrumMaster wird innerhalb des oben genannten Teams ernannt.

Folgeprojekte von \textit{Fantastic Feasts} sind derzeit nicht vorgesehen. Denkbar wäre jedoch ein Publishing des Spiels mit ständig laufendem Server und beliebig vielen Server-Instanzen – so könnte \textit{Fantastic Feasts} weltweit von zahlreichen Spielern gespielt werden.
    \section{Fachwissen}
    \begriff{Nutzer}
{Ein Mensch, der einen Rechner bedient und entweder den Client zum Spielen des Spiels oder zur Beobachtung einer Partie benutzt, oder den Team-Editor bedient. Jeder Benutzer hat einen Nutzernamen, mittels dem er von anderen Nutzern erkannt werden kann.}
{-}
{Spieler, Gast}
{Zur Beschreibung des Programmverlaufs}
{-}
{JägerMaister69}

\begriff{Spieler}
{Ein Nutzer, der das Computerspiel \glqq{}Fantastic Feasts\grqq{}  spielt.}
{Nutzer}
{-}
{Zur Beschreibung des Programmverlaufs}
{-}
{-}

\begriff{Gast}
{Ein Nutzer, der eine laufende Partie beobachtet}
{Nutzer}
{-}
{Zur Beschreibung des Programmverlaufs}
{-}
{-}

\begriff{Client}
{Das Computerprogramm, das mit einer grafischen Oberfläche ausgestattet ist und einem Nutzer erlaubt, eine Verbindung mit einem Server herzustellen und damit zu Kommunizieren}
{-}
{-}
{Zum Spielen des Spiels \glqq{}Fantastic Feasts\grqq{} }
{Der Begriff bezieht sich nicht auf den Menschen, der das Programm bedient}
{-}

\begriff{Server}
{Die zentrale Komponente, in dem die Spiellogik implementiert ist und die Programmbefehle abwickelt und mit dem sich Clients verbinden können, um eine Partie zu spielen oder zu beobachten. Die Kommunikation erfolgt mit JSON}
{-}
{-}
{Ist für die Kommunikation von Clients, für das Verwalten des Spielgeschehens, Ressourcenverwaltung und die Spiellogik verantwortlich.}
{-}
{-}

\begriff{Team-Editor}
{Ermöglicht einem Nutzer mit einer grafischen Oberfläche, ein eigenes Team zu erstellen und zu bearbeiten. Die Einstellungen werden danach als JSON-Datei gespeichert.}
{-}
{-}
{Zur Erstellung von nutzereigenen Teams.}
{-}
{-}

\begriff{KI-Client}
{Meldet sich beim Server wie ein normaler Client an und simuliert mit einer KI einen menschlichen Spieler. Hat keine grafische Oberfläche. Meldet sich mit dem Nutzernamen \glqq{}KI\grqq{}  ein.}
{-}
{-}
{Zum Spielen gegen einen Computergegner}
{-}
{-}

\begriff{KI}
{Definiert die Regeln, nach denen der KI-Client auf die durch den Server vermittelten Geschehen im Spiel reagiert.}
{-}
{-}
{Zum Spielen gegen einen Computergegner}
{Die KI ist die Logik, nach der der Computer das Spiel spielt und kein Programm}
{-}

\begriff{Spielfeld}
{Ein grafisch darstellbares Raster, auf dem sich die Spielfiguren bewegen}
{-}
{-}
{Dient als virtuelles Spielbrett mit klar definierten Abgrenzungen}
{Wird nicht Spielumgebung genannt um Verwechslung mit dem Client zu vermeiden}
{-}

\begriff{Zelle}
{Die kleinste Einheit des Spielfeldes, also ein Quadrat davon}
{-}
{Zentrumszelle, Torring, Hüterzonenzelle}
{Mögliche Standorte der Spielfiguren}
{Wird nicht Feld genannt, da das ein eher vager Begriff ist}
{-}

\begriff{Zentrum}
{Der 3x3 Zellen große Abschnitt in der Mitte des Spielfeldes}
{-}
{-}
{Summe aller Zentrumszellen}
{Ist das Mittelfeld im Lastenheft}
{-}

\begriff{Hüterzone}
{Die Bereiche am linken und rechten Rand des Spielfeldes, in dem sich die Torringe befinden}
{-}
{-}
{Summe aller kritischen Zellen und Torring}
{-}
{-}

\begriff{Torring}
{Die Zellen in die beide Teams den Quaffel bewegen wollen. Es wird zwischen eigenen und gegnerischen Torringen unterschieden.}
{Zelle}
{Eigener Torring, Gegnerischer Torring}
{Hauptquelle von Punkten}
{Torring im Lastenheft}
{}

\begriff{Zentrumszelle}
{Eine Zelle im Zentrum des Spielfeldes (siehe Zentrum)}
{Zelle}
{-}
{Startpunkt für Quaffel und Klatscher}
{-}
{-}

\begriff{Hüterzonenzelle}
{Eine Zelle in einem kritischen Bereich des Spielfeldes}
{Zelle}
{-}
{limitierendes Element für das Abliefern der Quaffel}
{-}
{-}

\begriff{Spielobjekt}
{Jedes Objekt, das sich auf dem Spielfeld befindet und darauf bewegt werden kann}
{-}
{Ball, Spielfigur}
{-}
{Nicht zu verwechseln mit Spielfigur}
{-}

\begriff{Ball}
{Ein Spielobjekt, das nicht direkt, nur indirekt von einem Spieler beeinflusst werden kann}
{Spielobjekt}
{Quaffel, Klatscher, Schnatz}
{Festpunkte zur Steuerung des Spielverlaufs}
{-}
{-}

\begriff{Spielfigur}
{Ein Spielobjekt, das von einem Spieler direkt gesteuert wird. Jede Spielfigur hat einen Namen, einen Besenrang, ein Geschlecht, ein Team und eine Rolle. Man unterscheidet außerdem zwischen eigenen und gegnerischen Spielfiguren.}
{Spielobjekt}
{Hüter, Sucher, Jäger, Treiber}
{Mitglieder eines Teams}
{Spieler im Lastenheft}
{Name: Luke Skywalker, Rolle: Hüter des Teams \glqq{}Jedi\grqq{}, Geschlecht: männlich, Besenrang: 5}

\begriff{Quaffel}
{Passives Objekt, mit dem Jäger und Hüter interagieren können und von ihnen nach Möglichkeit in ein gegnerisches Zielfeld befördert werden soll.}
{Ball}
{-}
{Zentrales Spielobjekt}
{-}
{-}

\begriff{Klatscher}
{Ball, der sich von selbst auf Spielfiguren zubewegt, die keine Treiber sind und diese betäuben können und von Treibern bewegt werden können.}
{Ball}
{-}
{Zusätzliches taktisches Spielelement}
{-}
{-}

\begriff{Schnatz}
{Ball, der von den Suchern gejagt wird und deren Fund die Partie beendet}
{Ball}
{-}
{Definiert Spielende}
{\glqq{}Schnatz\grqq{}  im Lastenheft}
{-}

\begriff{Partie}
{Ein einzelnes Spiel. Beginnt beim Platzieren der Figuren und endet mit dem Bestimmen des Gewinners.}
{-}
{-}
{Beschreibung des Spielablaufs}
{-}
{Spieler VodkaVodka98 spielt gegen Spieler LongEiländ}

\begriff{Hüter}
{Spielfigur, deren Aufgabe es ist, den Quaffel von den eigenen Zielfeldern fernzuhalten}
{Spielfigur}
{Eigener Hüter, Gegnerischer Hüter}
{Letzte Verteidigungslinie}
{-}
{Siehe \glqq{}Spielfigur\grqq{} }

\begriff{Sucher}
{Spielfigur, die den Schnatz jagt}
{Spielfigur}
{Eigener Sucher, Gegnerischer Sucher}
{Beendet die Partie}
{-}
{Darth Vader, gegnerischer Sucher, Besenrang 2}

\begriff{Jäger}
{Spielfigur, die den Quaffel in ein einen gegnerischen Torring befördern soll}
{Spielfigur}
{Eigener Jäger, Gegnerischer Jäger}
{Holt Punkte für das eigene Team}
{\glqq{}Jäger\grqq{}  im Lastenheft. Jäger beschreibt die Rolle der Spielfigur aber besser.}
{Han Solo, eigener Jäger, Besenrang 3}

\begriff{Treiber}
{Spielfigur, mit der der Spieler eigene Spielfiguren vor Klatschern schützt und gegnerische damit abschießen kann}
{Spielfigur}
{Eigener Treiber, Gegnerischer Treiber}
{Interagiert mit Klatschern}
{-}
{Boba Fett, gegnerischer Treiber, Besenrang 4}

\begriff{Geschlecht}
{Jede Spielfigur ist entweder männlich oder weiblich.}
{-}
{-}
{Team-Editierung}
{-}
{-}

\begriff{Team}
{Die Menge aller Spielfiguren auf dem Spielfeld, die von einem einzigen Spieler kontrolliert wird. Ein Team hat einen Namen, ein Motto, ein Logo, eine Hauptfarbe und eine Ersatzfarbe.}
{eigenes Team, gegnerisches Team}
{-}
{Beschreibung einer Partie}
{-}
{Galaktisches Imperium, Motto: \glqq{}Unbegrenzte MAAACHT!\grqq{}, [Todesstern als Logo], Hauptfarbe: Schwarz, Ersatzfarbe: Rot }

\begriff{Punkte}
{Der Spieler mit mehr Punkten am Ende der Partie gewinnt. Werden durch das Platzieren des Quaffel in einem gegnerischen Torring oder das Finden des Schnatzes erhalten.}
{-}
{-}
{Bestimmung des Gewinners}
{-}
{SchnapsNase hat 20 Punkte}

\begriff{Besetzen}
{Eine Spielfigur besetzt das Feld, auf dem sie sich befindet}
{-}
{-}
{Beschreibung des Spielgeschehens}
{Zwei Spielfiguren können sich nicht auf derselben Zelle befinden}
{Chewbacca besetzt Zelle 5:3}

\begriff{Besenrang}
{Jede Spielfigur hat einen Besenrang von 1 bis 5, der die Wahrscheinlichkeit bestimmt, dass sie noch einmal ziehen kann. Besenrang 1 ist der beste.}
{-}
{-}
{Unterscheidet Qualität der Spielfiguren.}
{Ersetzt die verschiedenen \glqq{}Besen\grqq{} aus dem Lastenheft mit einer Skala von 1 bis 5 zur besseren Übersicht.}
{Yoda hat Besenrang 1.}

\begriff{Aktion}
{Jede durch einen Spieler hervorgerufene Änderung der Spielsituation}
{-}
{Ziehen, Schießen, Schlagen, Einmischung, Übernahme}
{Weiterführung der Partie}
{-}
{Obi-Wan Kenobi zieht von Zelle 8:7 auf Zelle 9:7}

\begriff{Ziehen}
{Die Bewegung einer Spielfigur von einer Zelle auf eine andere durch direkten Befehl des Spielers}
{Aktion}
{-}
{Beschreibung des Spielverlaufs}
{Bezieht sich nicht auf erzwungene Bewegungen einer Spielfigur.}
{Obiwan Kenobi zieht von Zelle 8:7 auf Zelle 9:7}

\begriff{Befördern}
{Bewegen des Quaffel mittels einer Spielfigur}
{-}
{-}
{Bewegen des Quaffel, allgemeiner Begriff}
{Keine Aktion, da eventuell eine passive Folge, z.B. durch Ziehen}
{-}

\begriff{Schießen}
{Die Bewegung des Quaffel durch einen Hüter oder Jäger auf eine andere, entfernte Zelle ohne Bewegung der Spielfigur}
{Aktion}
{-}
{Bewegung des Quaffel um mehrere Felder}
{\glqq{}Werfen\grqq{}  im Lastenheft. Analog zum Schussvektor benannt.}
{Mace Windu schießt den Quaffel auf Zelle 10:4}

\begriff{Schlagen}
{Die erzwungene Bewegung eines Klatschers durch einen Treiber}
{Aktion}
{-}
{Interaktion mit Klatschern}
{\glqq{}Kloppen\grqq{}  im Lastenheft}
{R2-D2 schlägt einen Klatscher auf Zelle 5:10}

\begriff{Einmischung}
{Hilfsfähigkeiten, die nicht von Spielobjekten ausgehen. Werden von einem Spieler gesteuert. Bei jeder Benutzung besteht eine Chance, dass die verwendete Einmischung bis zum Ende der Partie für den jeweiligen Spieler vom Schiedsrichter deaktiviert werden.}
{Aktion}
{Teleportation, Fernangriff, Impuls, Schnatzjagd}
{Zusätzliche taktische Element}
{Ersetzt die \glqq{}Fans\grqq{}  aus dem Lastenheft}
{Lando Calrissian wird auf Zelle 6:6 teleportiert}

\begriff{Teleportation}
{Einmischung, die eine Spielfigur auf eine zufällige Zelle teleportiert}
{Einmischung}
{-}
{-}
{Ersetzt \glqq{}Elfen\grqq{}  aus Lastenheft}
{Siehe \glqq{}Einmischungen\grqq{} }

\begriff{Fernangriff}
{Trifft eine gegnerische Spielfigur. Ziel verliert gegebenenfalls den Quaffel und wird auf eine zufällige benachbarte, freie Zelle bewegt.}
{Einmischung}
{-}
{-}
{Statt \glqq{}Kobolde\grqq{}  im Lastenheft}
{Jango Fett wird von Fernangriff auf Zelle 5:6 gestoßen}

\begriff{Impuls}
{Wenn eine Spielfigur den Quaffel hält, wird sie bei Benutzung verloren}
{Einmischung}
{-}
{-}
{Statt \glqq{}Trolle\grqq{}  im Lastenheft}
{C-3PO verliert wegen eines Impuls den Quaffel}

\begriff{Schnatzstoß}
{Bewegt den Schnatz in eine zufällige Richtung um ein Feld}
{Einmischung}
{-}
{-}
{\glqq{}Schnatzschnappen\grqq{}  im Lastenheft}
{Ein Schnatzstoß treibt den Schnatz auf Zelle 4:12}

\begriff{Entfernung}
{Eine Entfernung zwischen zwei Zellen ist die minimale Anzahl von Zügen, in denen eine Spielfigur von der einen auf die andere ziehen kann.}
{-}
{-}
{Spielfeldgeometrie}
{-}
{-}

\begriff{Schussvektor}
{Pfeil vom Mittelpunkt einer Zelle zum Mittelpunkt einer anderen}
{-}
{Torschussvektor}
{Spielfeldgeometrie}
{-}
{-}

\begriff{Torschussvektor}
{Schussvektor zu einem Schuss, der möglicherweise in einem Torschuss resultiert. (Ein Schussvektor, der die linke oder rechte Seite eines Torrings schneidet.)}
{Schussvektor}
{-}
{Punkte sammeln}
{-}
{-}

\begriff{Torschuss}
{Ein Jäger schießt den Quaffel in einen Torring und holt damit Punkte für sei Team}
{-}
{-}
{Punkte sammeln}
{Nur erfolgreiche Schüsse auf das Tor werden als Torschüsse bezeichnet.}
{Darth Sidious schießt den Quaffel in ein eigenes Tor}

\begriff{Zugphase}
{Phase, in der eine Spielfigur Aktionen durchführt. Beginnt, sobald der Spieler die Möglichkeit hat, die jeweilige Spielfigur zu steuern und endet, sobald er ihr den letzten Befehl für diesen Zug gegeben hat. Ein Zug enthält mehrere Zugphasen.}
{-}
{-}
{Zeiteinteilung}
{}
{Leia Organa ist dran.}

\begriff{Zug}
{Von der ersten Aktion eines Spielers bis zur ersten Aktion des Gegners}
{-}
{-}
{Zeiteinteilung}
{Nicht die Zugphase einer Spielfigur}
{Bierdurst69 ist am Zug}

\begriff{Endphase}
{Letzter Teil eines Zuges. Der Spieler kann darin Einmischungen vornehmen.}
{-}
{-}
{Zeiteinteilung}
{-}
{-}

\begriff{Verlieren}
{Der Quaffel wird auf eine zufällige Zelle bewegt, die an die Zelle angrenzt, auf der sich die Spielfigur, die sie derzeit hält befindet.}
{-}
{-}
{Spielablauf}
{\glqq{}Vertändeln\grqq{}  im Lastenheft}
{Jar Jar verliert den Quaffel}

\begriff{Halten}
{Ein Jäger oder Hüter kann den Quaffel halten. Ist das der Fall, bewegt sich der Quaffel auf die Zelle, auf die die Spielfigur zieht.}
{-}
{-}
{Beschreibung des Spielgeschehens}
{-}
{-}

\begriff{Übernahme}
{Ein Jäger neben einer gegnerischen Spielfigur, die den Quaffel hält, kann diesen mit einer bestimmten Wahrscheinlichkeit übernehmen und hält sie anschließend selbst.}
{Aktion}
{-}
{Aggressives Spielmanöver}
{-}
{Darth Vader übernimmt den Quaffel von Anakin Skywalker}

\begriff{Betäubt}
{Eine betäubte Spielfigur kann in seiner nächsten Zugphase keine Aktion durchführen}
{-}
{-}
{Wirkung der Klatscher}
{\glqq{}Ausgeknockt\grqq{}  im Lastenheft}
{Jango Fett ist betäubt}

\begriff{Foul}
{Handlung, wegen der eine Spielfigur vorübergehend vom Spielfeld entfernt werden kann.}
{-}
{Torring Blockieren, Stürmen, Großoffensive, Rammen, Schnatz Blockieren}
{Taktische Elemente}
{-}
{Qui-Gon Jinn blockiert den Schnatz}

\begriff{Torring Blockieren}
{Eine eigene Spielfigur besetzt einen eigenen Torring, was eine Torschuss verhindert.}
{Foul}
{-}
{Taktik}
{\glqq{}Flackern\grqq{}  im Pflichtenheft}
{-}

\begriff{Stürmen}
{Ein Jäger, der den Quaffel hält, zieht auf einen gegnerischen Torring, was das Abliefern garantiert.}
{Foul}
{-}
{Taktik}
{\glqq{}Nachtarocken\grqq{}  im Lastenheft}
{Han Solo stürmt mittleren gegnerischen Torring}

\begriff{Großoffensive}
{Ein eigener Jäger betritt eine gegnerische Hüterzonenzelle während ein anderer eigener Jäger sich auf einer anderen befindet.}
{Foul}
{-}
{Taktik}
{\glqq{}Stutschen\grqq{}  im Lastenheft}
{Lando Calrissia schließt sich Chewbacca in einer Großoffensive an}

\begriff{Rammen}
{Eine eigene Spielfigur zieht auf eine Zelle, die von einer gegnerischen Spielfigur besetzt wird. Dadurch wird die gegnerische Spielfigur auf eine benachbarte Zelle bewegt und verliert den Quaffel}
{Foul}
{-}
{Taktik}
{\glqq{}Keilen\grqq{}  im Lastenheft}
{Boba Fett rammt Jar Jar}

\begriff{Schnatz blockieren}
{Eine Spielfigur, die kein Sucher ist, besetzt die Zelle, auf der sich der Schnatz befindet.}
{Foul}
{-}
{Taktik}
{\glqq{}Schnatzeln\grqq{}  im Lastenheft}
{Darth Maul blockiert den Schnatz}

\begriff{Schiedsrichter}
{Entfernt mit bestimmter Wahrscheinlichkeit eine Spielfigur, die ein Foul ausführt vom Spielfeld bis ein Torschuss erfolgt und deaktiviert permanent eine Einmischung für den Rest der Partie.}
{-}
{-}
{Taktik}
{\glqq{}Schiedsrichter\grqq{}  im Lastenheft}
{Sheev Palpatine wurde vom Schiedsrichter vom Spielfeld entfernt}

\begriff{Disqualifikation}
{Tritt ein wenn fünf Spielfiguren eines Spielers gleichzeitig durch den Schiedsrichter aus dem Spiel entfernt sind. Führt zur Niederlage des Spielers.}
{-}
{-}
{Erhöhtes Risiko}
{-}
{CubaLibre wurde disqualifiziert. CaptainCola gewinnt die Partie.}

    \includepdf[pagecommand={\section{Domänenmodell}}]{domaenenmodell.pdf}
    \section{Anforderungsdefinition}
    %TODO

    \subsection{Akteure}

    \fakt	{Anwendungsbenutzer}
            {Menschlicher Benutzer, der die Anwendungen bedient.}
            {<Rolle>}

    \fakt	{Spielender Benutzer}
            {Anwendungsbenutzer, der mit der Client-Anwendung am tatsächlichen Spielgeschehen teilnimmt.}
            {<Rolle>}

    \fakt	{Beobachtender Benutzer}
            {<Beschreibung>}
            {<Rolle>}

    \fakt	{Systemadministrator}
            {Anwendungsnutzer mit Zusatzqualifikation um Server zu verwalten.}
            {Der Systemadministrator hat zugriff auf die Server Anwendung. Er ist dafür verantwortlich eine Instanz der Server Anwendung zu starten und zu betreuen. Zudem hat er Zugriff auf die Partie-Konfiguration und kann diese bei bedarf verändern.}

    \fakt	{KI}
            {<Beschreibung>}
            {<Rolle>}

    \fakt	{Kunde}
            {SoPra-Tutor}
            {<Rolle>}

    \fakt	{Client [Komponente]}
            {<Beschreibung>}
            {<Rolle>}

    \fakt	{KI-Client [Komponente]}
            {<Beschreibung>}
            {<Rolle>}

    \fakt	{Server [Komponente]}
            {<Beschreibung>}
            {<Rolle>}

    \fakt	{Quidditchteam-Editor [Komponente]}
            {<Beschreibung>}
            {<Rolle>}

    \fakt 	{Team}
            {<Beschreibung>}
            {<Rolle>}

    \fakt 	{Schiedsrichter}
            {<Beschreibung>}
            {<Rolle>}


    \subsection{Funktionale Anforderungen (allgemein)}

    \fanf	{Quidditch-Spielfeld}
            {<Beschreibung>}
            {Diese Anforderung geht aus den im Lastenheft zu Verfügung gestellten Spielregeln für das Spiel \textit{Fantastic Feasts} hervor.}
            {<Abhängigkeiten>}
            {<Prio>}
            {<Akteure>}

    \fanf	{Mittelkreis}
            {<Beschreibung>}
            {Diese Anforderung geht aus den im Lastenheft zu Verfügung gestellten Spielregeln für das Spiel \textit{Fantastic Feasts} hervor.}
            {<Abhängigkeiten>}
            {<Prio>}
            {<Akteure>}

    \fanf	{Mittelfeld}
            {<Beschreibung>}
            {Diese Anforderung geht aus den im Lastenheft zu Verfügung gestellten Spielregeln für das Spiel \textit{Fantastic Feasts} hervor.}
            {<Abhängigkeiten>}
            {<Prio>}
            {<Akteure>}

    \fanf	{Hüterzone}
            {<Beschreibung>}
            {<Begründung>}
            {<Abhängigkeiten>}
            {<Prio>}
            {<Akteure>}

    \fanf	{Torring}
            {<Beschreibung>}
            {Diese Anforderung geht aus den im Lastenheft zu Verfügung gestellten Spielregeln für das Spiel \textit{Fantastic Feasts} hervor.}
            {<Abhängigkeiten>}
            {<Prio>}
            {<Akteure>}

    \fanf	{Schussvektor}
            {<Beschreibung>}
            {Diese Anforderung geht aus den im Lastenheft zu Verfügung gestellten Spielregeln für das Spiel \textit{Fantastic Feasts} hervor.}
            {<Abhängigkeiten>}
            {<Prio>}
            {<Akteure>}

    \fanf	{Bälle}
            {<Beschreibung>}
            {Diese Anforderung geht aus den im Lastenheft zu Verfügung gestellten Spielregeln für das Spiel \textit{Fantastic Feasts} hervor.}
            {<Abhängigkeiten>}
            {<Prio>}
            {<Akteure>}

    \fanf	{Quaffel [Ball]}
            {<Beschreibung>}
            {Diese Anforderung geht aus den im Lastenheft zu Verfügung gestellten Spielregeln für das Spiel \textit{Fantastic Feasts} hervor.}
            {<Abhängigkeiten>}
            {<Prio>}
            {<Akteure>}

    \fanf	{Klatscher [Ball]}
            {<Beschreibung>}
            {Diese Anforderung geht aus den im Lastenheft zu Verfügung gestellten Spielregeln für das Spiel \textit{Fantastic Feasts} hervor.}
            {<Abhängigkeiten>}
            {<Prio>}
            {<Akteure>}

    \fanf	{Goldener Schnatz [Ball]}
            {<Beschreibung>}
            {Diese Anforderung geht aus den im Lastenheft zu Verfügung gestellten Spielregeln für das Spiel \textit{Fantastic Feasts} hervor.}
            {<Abhängigkeiten>}
            {<Prio>}
            {<Akteure>}

    \fanf	{Besen}
            {<Beschreibung>}
            {Diese Anforderung geht aus den im Lastenheft zu Verfügung gestellten Spielregeln für das Spiel \textit{Fantastic Feasts} hervor.}
            {<Abhängigkeiten>}
            {<Prio>}
            {<Akteure>}

    \fanf	{Zauberfauch [Besen]}
            {<Beschreibung>}
            {Diese Anforderung geht aus den im Lastenheft zu Verfügung gestellten Spielregeln für das Spiel \textit{Fantastic Feasts} hervor.}
            {<Abhängigkeiten>}
            {<Prio>}
            {<Akteure>}

    \fanf	{Sauberwisch 11 [Besen]}
            {<Beschreibung>}
            {Diese Anforderung geht aus den im Lastenheft zu Verfügung gestellten Spielregeln für das Spiel \textit{Fantastic Feasts} hervor.}
            {<Abhängigkeiten>}
            {<Prio>}
            {<Akteure>}

    \fanf	{Komet 2-60 [Besen]}
            {<Beschreibung>}
            {Diese Anforderung geht aus den im Lastenheft zu Verfügung gestellten Spielregeln für das Spiel \textit{Fantastic Feasts} hervor.}
            {<Abhängigkeiten>}
            {<Prio>}
            {<Akteure>}

    \fanf	{Nimbus 2001  [Besen]}
            {<Beschreibung>}
            {Diese Anforderung geht aus den im Lastenheft zu Verfügung gestellten Spielregeln für das Spiel \textit{Fantastic Feasts} hervor.}
            {<Abhängigkeiten>}
            {<Prio>}
            {<Akteure>}

    \fanf	{Feuerblitz [Besen]}
            {<Beschreibung>}
            {Diese Anforderung geht aus den im Lastenheft zu Verfügung gestellten Spielregeln für das Spiel \textit{Fantastic Feasts} hervor.}
            {<Abhängigkeiten>}
            {<Prio>}
            {<Akteure>}

    \fanf	{Spieler}
            {<Beschreibung>}
            {Diese Anforderung geht aus den im Lastenheft zu Verfügung gestellten Spielregeln für das Spiel \textit{Fantastic Feasts} hervor.}
            {<Abhängigkeiten>}
            {<Prio>}
            {<Akteure>}	

    \fanf	{Jäger [Spielertyp]}
            {<Beschreibung>}
            {Diese Anforderung geht aus den im Lastenheft zu Verfügung gestellten Spielregeln für das Spiel \textit{Fantastic Feasts} hervor.}
            {<Abhängigkeiten>}
            {<Prio>}
            {<Akteure>}

    \fanf	{Treiber [Spielertyp]}
            {<Beschreibung>}
            {Diese Anforderung geht aus den im Lastenheft zu Verfügung gestellten Spielregeln für das Spiel \textit{Fantastic Feasts} hervor.}
            {<Abhängigkeiten>}
            {<Prio>}
            {<Akteure>}

    \fanf	{Hüter [Spielertyp]}
            {<Beschreibung>}
            {Diese Anforderung geht aus den im Lastenheft zu Verfügung gestellten Spielregeln für das Spiel \textit{Fantastic Feasts} hervor.}
            {<Abhängigkeiten>}
            {<Prio>}
            {<Akteure>}

    \fanf	{[Spielertyp] Sucher}
            {<Beschreibung>}
            {Diese Anforderung geht aus den im Lastenheft zu Verfügung gestellten Spielregeln für das Spiel \textit{Fantastic Feasts} hervor.}
            {<Abhängigkeiten>}
            {<Prio>}
            {<Akteure>}

    \fanf	{Fans]}
            {<Beschreibung>}
            {Diese Anforderung geht aus den im Lastenheft zu Verfügung gestellten Spielregeln für das Spiel \textit{Fantastic Feasts} hervor.}
            {<Abhängigkeiten>}
            {<Prio>}
            {<Akteure>}	

    \fanf	{Elfen [Fantyp]}
            {<Beschreibung>}
            {Diese Anforderung geht aus den im Lastenheft zu Verfügung gestellten Spielregeln für das Spiel \textit{Fantastic Feasts} hervor.}
            {<Abhängigkeiten>}
            {<Prio>}
            {<Akteure>}

    \fanf	{Kobolde [Fantyp]}
            {<Beschreibung>}
            {Diese Anforderung geht aus den im Lastenheft zu Verfügung gestellten Spielregeln für das Spiel \textit{Fantastic Feasts} hervor.}
            {<Abhängigkeiten>}
            {<Prio>}
            {<Akteure>}

    \fanf	{Trolle [Fantyp]}
            {<Beschreibung>}
            {Diese Anforderung geht aus den im Lastenheft zu Verfügung gestellten Spielregeln für das Spiel \textit{Fantastic Feasts} hervor.}
            {<Abhängigkeiten>}
            {<Prio>}
            {<Akteure>}

    \fanf	{Niffler [Fantyp]}
            {<Beschreibung>}
            {Diese Anforderung geht aus den im Lastenheft zu Verfügung gestellten Spielregeln für das Spiel \textit{Fantastic Feasts} hervor.}
            {<Abhängigkeiten>}
            {<Prio>}
            {<Akteure>}

    \fanf	{Foul}
            {<Beschreibung>}
            {Diese Anforderung geht aus den im Lastenheft zu Verfügung gestellten Spielregeln für das Spiel \textit{Fantastic Feasts} hervor.}
            {<Abhängigkeiten>}
            {<Prio>}
            {<Akteure>}

    \fanf	{Flacken [Foul]}
            {<Beschreibung>}
            {Diese Anforderung geht aus den im Lastenheft zu Verfügung gestellten Spielregeln für das Spiel \textit{Fantastic Feasts} hervor.}
            {<Abhängigkeiten>}
            {<Prio>}
            {<Akteure>}

    \fanf	{Nachtarocken [Foul]}
            {<Beschreibung>}
            {Diese Anforderung geht aus den im Lastenheft zu Verfügung gestellten Spielregeln für das Spiel \textit{Fantastic Feasts} hervor.}
            {<Abhängigkeiten>}
            {<Prio>}
            {<Akteure>}

    \fanf	{Stutschen [Foul]}
            {<Beschreibung>}
            {Diese Anforderung geht aus den im Lastenheft zu Verfügung gestellten Spielregeln für das Spiel \textit{Fantastic Feasts} hervor.}
            {<Abhängigkeiten>}
            {<Prio>}
            {<Akteure>}

    \fanf	{Keilen [Foul]}
            {<Beschreibung>}
            {Diese Anforderung geht aus den im Lastenheft zu Verfügung gestellten Spielregeln für das Spiel \textit{Fantastic Feasts} hervor.}
            {<Abhängigkeiten>}
            {<Prio>}
            {<Akteure>}

    \fanf	{Schnaltzeln [Foul]}
            {<Beschreibung>}
            {Diese Anforderung geht aus den im Lastenheft zu Verfügung gestellten Spielregeln für das Spiel \textit{Fantastic Feasts} hervor.}
            {<Abhängigkeiten>}
            {<Prio>}
            {<Akteure>}

    \fanf	{Wurf mit dem Quaffel}
            {<Beschreibung>}
            {Diese Anforderung geht aus den im Lastenheft zu Verfügung gestellten Spielregeln für das Spiel \textit{Fantastic Feasts} hervor.}
            {<Abhängigkeiten>}
            {<Prio>}
            {<Akteure>}

    \fanf	{Quaffel Abfangen}
            {<Beschreibung>}
            {Diese Anforderung geht aus den im Lastenheft zu Verfügung gestellten Spielregeln für das Spiel \textit{Fantastic Feasts} hervor.}
            {<Abhängigkeiten>}
            {<Prio>}
            {<Akteure>}

    \fanf	{Torschuss}
            {<Beschreibung>}
            {Diese Anforderung geht aus den im Lastenheft zu Verfügung gestellten Spielregeln für das Spiel \textit{Fantastic Feasts} hervor.}
            {<Abhängigkeiten>}
            {<Prio>}
            {<Akteure>}

    \fanf	{Klatscher kloppen}
            {<Beschreibung>}
            {Diese Anforderung geht aus den im Lastenheft zu Verfügung gestellten Spielregeln für das Spiel \textit{Fantastic Feasts} hervor.}
            {<Abhängigkeiten>}
            {<Prio>}
            {<Akteure>}

    \fanf	{Partie-Konfiguration}
            {<Beschreibung>}
            {<Begründung>}
            {<Abhängigkeiten>}
            {<Prio>}
            {<Akteure>}

    \fanf	{Quidditchteam-Konfiguration}
            {<Beschreibung>}
            {<Begründung>}
            {<Abhängigkeiten>}
            {<Prio>}
            {<Akteure>}	

    \fanf	{Netzwerkinterface}
            {<Beschreibung>}
            {<Begründung>}
            {<Abhängigkeiten>}
            {<Prio>}
            {<Akteure>}

    \fanf	{Runde}
            {<Beschreibung>}
            {Diese Anforderung geht aus den im Lastenheft zu Verfügung gestellten Spielregeln für das Spiel \textit{Fantastic Feasts} hervor.}
            {<Abhängigkeiten>}
            {<Prio>}
            {<Akteure>}

    \fanf	{Ballphase}
            {<Beschreibung>}
            {Diese Anforderung geht aus den im Lastenheft zu Verfügung gestellten Spielregeln für das Spiel \textit{Fantastic Feasts} hervor.}
            {<Abhängigkeiten>}
            {<Prio>}
            {<Akteure>}

    \fanf	{Spielerphase}
            {<Beschreibung>}
            {Diese Anforderung geht aus den im Lastenheft zu Verfügung gestellten Spielregeln für das Spiel \textit{Fantastic Feasts} hervor.}
            {<Abhängigkeiten>}
            {<Prio>}
            {<Akteure>}

    \fanf	{Fanphase}
            {<Beschreibung>}
            {Diese Anforderung geht aus den im Lastenheft zu Verfügung gestellten Spielregeln für das Spiel \textit{Fantastic Feasts} hervor.}
            {<Abhängigkeiten>}
            {<Prio>}
            {<Akteure>}

    \fanf	{Zufallsgenerator}
            {<Beschreibung>}
            {<Begründung>}
            {<Abhängigkeiten>}
            {<Prio>}
            {<Akteure>}

    \fanf	{Spielende}
            {<Beschreibung>}
            {<Begründung>}
            {<Abhängigkeiten>}
            {<Prio>}
            {<Akteure>}

    \fanf	{Überlängenbehnadlung}
            {<Beschreibung>}
            {<Begründung>}
            {<Abhängigkeiten>}
            {<Prio>}
            {<Akteure>}

    \fanf	{Spiellogik}
            {<Beschreibung>}
            {<Begründung>}
            {<Abhängigkeiten>}
            {<Prio>}
            {<Akteure>}

    \fanf	{Log-Datei}
            {<Beschreibung>}
            {<Begründung>}
            {<Abhängigkeiten>}
            {<Prio>}
            {<Akteure>}

    \fanf	{<Titel>}
            {<Beschreibung>}
            {<Begründung>}
            {<Abhängigkeiten>}
            {<Prio>}
            {<Akteure>}

    \fanf	{<Titel>}
            {<Beschreibung>}
            {<Begründung>}
            {<Abhängigkeiten>}
            {<Prio>}
            {<Akteure>}

    \fanf	{<Titel>}
            {<Beschreibung>}
            {<Begründung>}
            {<Abhängigkeiten>}
            {<Prio>}
            {<Akteure>}

    \fanf	{<Titel>}
            {<Beschreibung>}
            {<Begründung>}
            {<Abhängigkeiten>}
            {<Prio>}
            {<Akteure>}

    \fanf	{<Titel>}
            {<Beschreibung>}
            {<Begründung>}
            {<Abhängigkeiten>}
            {<Prio>}
            {<Akteure>}


    \subsection{Funktionale Anforderungen (Server)}

    \subsection{Funktionale Anforderungen (Client)}

    \fanf	{Hauptmenü [Client-View]}
            {<Beschreibung>}
            {<Begründung>}
            {<Abhängigkeiten>}
            {<Prio>}
            {<Akteure>}

    \fanf	{Connect to Game [Client-View]}
            {<Beschreibung>}
            {<Begründung>}
            {<Abhängigkeiten>}
            {<Prio>}
            {<Akteure>}

    \fanf	{End of Game [Client-View]}
            {<Beschreibung>}
            {<Begründung>}
            {<Abhängigkeiten>}
            {<Prio>}
            {<Akteure>}

    \fanf	{Import Team Config [Client-View]}
            {<Beschreibung>}
            {<Begründung>}
            {<Abhängigkeiten>}
            {<Prio>}
            {<Akteure>}

    \fanf	{Game Play [Client-View]}
            {<Beschreibung>}
            {<Begründung>}
            {<Abhängigkeiten>}
            {<Prio>}
            {<Akteure>}	

    \fanf	{Hilfe [Client-View]}
            {<Beschreibung>}
            {<Begründung>}
            {<Abhängigkeiten>}
            {<Prio>}
            {<Akteure>}

    \fanf	{Beobachter [Client-View]}
            {<Beschreibung>}
            {<Begründung>}
            {<Abhängigkeiten>}
            {<Prio>}
            {<Akteure>}

    \fanf	{Game Rendering Engine}
            {<Beschreibung>}
            {<Begründung>}
            {<Abhängigkeiten>}
            {<Prio>}
            {<Akteure>}

    \fanf	{Input Handler [Client]}
            {<Beschreibung>}
            {<Begründung>}
            {<Abhängigkeiten>}
            {<Prio>}
            {<Akteure>}

    \fanf	{Input Validierung [Client]}
            {<Beschreibung>}
            {<Begründung>}
            {<Abhängigkeiten>}
            {<Prio>}
            {<Akteure>}

    \fanf	{Hotkey [Client]}
            {<Beschreibung>}
            {<Begründung>}
            {<Abhängigkeiten>}
            {--}
            {<Akteure>}

    \fanf	{Pausieren [Client]}
            {<Beschreibung>}
            {<Begründung>}
            {<Abhängigkeiten>}
            {--}
            {<Akteure>}


    \subsection{Funktionale Anforderungen (Quidditchteam-Editor)}

    \fanf	{Edit Team [Team-Editor-View]}
            {<Beschreibung>}
            {<Begründung>}
            {<Abhängigkeiten>}
            {<Prio>}
            {<Akteure>}

    \fanf	{Save Team [Team-Editor-View]}
            {<Beschreibung>}
            {<Begründung>}
            {<Abhängigkeiten>}
            {<Prio>}
            {<Akteure>}


    \subsection{Nicht funktionale Anforderungen}

    \qanf 	{Plattformunabhängigkeit}
            {Der Spielclient soll auf mindestens einer gängigen Computerbetriebssystem-Plattform (z.B. Linux, Windows) uneingeschränkt benutzbar sein. Des weiteren soll die Server- und die KI-Komponente auf mindestens zwei gängigen Computerbetriebssystem-Plattformen (z.B. Linux, Windows) uneingeschränkt benutzbar sein.}
            {<Begründung>}
            {<Abhängigkeiten>}
            {<Prio>}
            {<Akteure>}

    \qanf 	{Version-Controlling}
            {Verwaltung und Version-Controlling des Quellcodes mit Hilfe eines Git basierten Version-Controlling Tool.}
            {<Begründung>}
            {-}
            {<Prio>}
            {<Akteure>}

    \qanf 	{Continuous Integration}
            {<Beschreibung>}
            {<Begründung>}
            {-}
            {<Prio>}
            {<Akteure>}

    \qanf 	{Statische Codeanalyse}
            {Mit Hilfe des Tools 'SonarQube' bzw. 'SonarCloud' soll der gesamt Quellcode einer statischen Analyse unterzogen werden. Dabei darf die technische Codequalität von diesen Tool nicht schlechter als 'B' bewertet werden.}
            {<Begründung>}
            {-}
            {<Prio>}
            {<Akteure>}

    \qanf 	{Qualitätssicherung}
            {80 Prozent des Quellcodes der einzelnen Komponenten soll durch automatische Test auf Korrektheit geprüft werden}
            {<Begründung>}
            {-}
            {<Prio>}
            {<Akteure>}

    \qanf 	{Docker Container}
            {<Beschreibung>}
            {<Begründung>}
            {-}
            {<Prio>}
            {<Akteure>}

    \qanf 	{Dokumentation}
            {<Beschreibung>}
            {<Begründung>}
            {-}
            {<Prio>}
            {<Akteure>}

    \qanf 	{Benutzerhandbuch}
            {<Beschreibung>}
            {<Begründung>}
            {-}
            {<Prio>}
            {<Akteure>}

    \qanf 	{Anwendungssprache}
            {<Beschreibung>}
            {<Begründung>}
            {-}
            {<Prio>}
            {<Akteure>}

    \qanf 	{Programmiersprache}
            {Implementierungssprache}
            {<Begründung>}
            {-}
            {<Prio>}
            {<Akteure>}

    \qanf 	{Dokumentationssprache}
            {<Beschreibung>}
            {<Begründung>}
            {-}
            {<Prio>}
            {<Akteure>}

    \qanf 	{Format für Konfigurationsdateien}
            {<Beschreibung>}
            {<Begründung>}
            {-}
            {<Prio>}
            {<Akteure>}

    \qanf 	{Netzwerkkommunikation}
            {<Beschreibung>}
            {<Begründung>}
            {-}
            {<Prio>}
            {<Akteure>}

    \qanf 	{Rundenspiel}
            {<Beschreibung>}
            {<Begründung>}
            {-}
            {<Prio>}
            {<Akteure>}

    \qanf 	{<Titel>}
            {<Beschreibung>}
            {<Begründung>}
            {-}
            {<Prio>}
            {<Akteure>}

    \qanf 	{<Titel>}
            {<Beschreibung>}
            {<Begründung>}
            {-}
            {<Prio>}
            {<Akteure>}

    \qanf 	{<Titel>}
            {<Beschreibung>}
            {<Begründung>}
            {-}
            {<Prio>}
            {<Akteure>}

    \qanf 	{<Titel>}
            {<Beschreibung>}
            {<Begründung>}
            {-}
            {<Prio>}
            {<Akteure>}

    \qanf 	{<Titel>}
            {<Beschreibung>}
            {<Begründung>}
            {-}
            {<Prio>}
            {<Akteure>}

    \qanf 	{<Titel>}
            {<Beschreibung>}
            {<Begründung>}
            {-}
            {<Prio>}
            {<Akteure>}


\end{document}
